\section{\tr{Dados}{Data}}
\label{sec:dados}

\pt{%
O estudo usa apenas dados que a empresa já coleta na rotina; não há coleta própria. Os números abaixo são ordens de grandeza, de um levantamento preliminar feito pelo autor, e serão atualizados no início do projeto.
}
\en{%
The study uses only data the company already collects routinely; there is no new data collection. The figures below are orders of magnitude from a preliminary survey by the author and will be updated at the start of the project.
}

\begin{table}[h]
\centering
\small
\caption{\tr{Fontes de dados previstas}{Planned data sources}}
\label{tab:fontes}
\begin{tabularx}{\textwidth}{>{\raggedright\arraybackslash}p{3.2cm} >{\raggedright\arraybackslash}X >{\raggedright\arraybackslash}p{3.6cm}}
\toprule
\tr{Fonte}{Source} & \tr{Conteúdo}{Content} & \tr{Escala aproximada}{Approximate scale} \\
\midrule
\tr{Relatos de acidentes}{Accident reports} &
\tr{Data, tipo de evento, unidade, número do veículo, gravidade, descrição}%
   {Date, event type, site, vehicle number, severity, description} &
\tr{Centenas de relatos em pouco mais de um ano}{Hundreds of reports over a little more than a year} \\
\tr{Telemetria veicular}{Vehicle telematics} &
\tr{Viagens (distância, duração), eventos de condução, colisões detectadas, unidade do veículo}%
   {Trips (distance, duration), driving events, detected collisions, vehicle site} &
\tr{Cerca de 300 veículos; mais de um milhão de viagens}{About 300 vehicles; over one million trips} \\
\tr{Horas trabalhadas}{Hours worked} &
\tr{Horas por unidade e mês}{Hours per site and month} &
\tr{Dezenas de unidades}{Dozens of sites} \\
\tr{Manutenção}{Maintenance} &
\tr{Ordens de reparo por veículo, com motivo (dano, acidente, preventiva)}%
   {Repair orders per vehicle, with reason (damage, accident, preventive)} &
\tr{Centenas de ordens por dano}{Hundreds of damage orders} \\
\bottomrule
\end{tabularx}
\end{table}

\pt{%
Um levantamento preliminar indicou que a ligação entre as fontes é viável: o número do veículo informado nos relatos coincide com o identificador da telemetria na maior parte dos casos, e o cadastro da telemetria associa cada veículo a uma unidade com o mesmo código usado no sistema de folha. O período comum às três fontes começa em 2025 \preencher{atualizar período e cobertura no início do projeto}.
}
\en{%
A preliminary survey indicated that linking the sources is feasible: the vehicle number given in the reports matches the telematics identifier in most cases, and the telematics registry assigns each vehicle to a site using the same code as the payroll system. The period common to the three sources starts in 2025 \preencher{update period and coverage at the start of the project}.
}

\pt{%
Os problemas de qualidade já conhecidos fazem parte do trabalho: relatos de teste misturados aos reais no período de implantação do aplicativo, unidade e veículo digitados em texto livre em parte dos relatos, viagens com registro de fim inválido na telemetria e mudanças de veículos entre unidades ao longo do tempo. Cada regra de limpeza será documentada e seu efeito sobre os resultados, testado em análise de sensibilidade.
}
\en{%
Known data quality issues are part of the work: test reports mixed with real ones during the application roll-out, site and vehicle typed as free text in some reports, trips with invalid end records in the telematics, and vehicles moving between sites over time. Every cleaning rule will be documented and its effect on the results tested in a sensitivity analysis.
}

\subsection*{\tr{Descrição preliminar}{Preliminary description}}

\pt{%
Uma análise descritiva agregada dos dados da operação, só com valores relativos (Figura~\ref{fig:desc-pp}), já mostra quatro fatos que orientam o desenho. (a)~As horas de viagem por veículo-mês são parecidas nos três tipos de uso, mas as milhas, não: o spotter anda a cerca de \descVelSpotter{} da velocidade do veículo local, o que justifica o tempo como exposição. (b)~Os acidentes parecem concentrados em poucos veículos (Gini de \descGiniObs), mas quase toda essa concentração é a esperada por acaso e exposição (\descGiniPoisson{} sob Poisson); concentração não é, por si só, evidência de risco maior. (c)~A hora do acidente está registrada em \descComHora\% dos relatos, e os acidentes se distribuem pelos três turnos, o que torna possíveis análises por turno. (d)~Relatos, ordens de manutenção por dano e colisões detectadas variam de forma diferente ao longo do ano, sinal de que captam universos diferentes. A dissertação terá um capítulo descritivo completo, com a distribuição de cada métrica usada nos modelos, seguido dos capítulos preditivo e prescritivo (Seção~\ref{sec:metodologia}).
}
\en{%
An aggregated descriptive analysis of the operational data, with relative values only (Figure~\ref{fig:desc-pp}), already shows four facts that shape the design. (a)~Trip hours per vehicle-month are similar across the three use types, but miles are not: the spotter moves at about \descVelSpotter{} of the local vehicle's speed, which justifies time as exposure. (b)~Accidents look concentrated in a few vehicles (Gini of \descGiniObs), but almost all of this concentration is what chance and exposure produce (\descGiniPoisson{} under Poisson); concentration is not, in itself, evidence of higher risk. (c)~The accident time is recorded in \descComHora\% of reports, and accidents are spread over the three shifts, which makes shift-level analyses possible. (d)~Reports, damage repair orders and detected collisions vary differently over the year, a sign that they capture different universes. The dissertation will have a full descriptive chapter, with the distribution of each metric used in the models, followed by the predictive and prescriptive chapters (Section~\ref{sec:metodologia}).
}

\begin{figure}[!htbp]
\centering
\includegraphics[width=\textwidth]{\arqfig{fig_desc_pp}}
\caption{\tr%
  {Descrição preliminar dos dados da operação (só valores relativos). (a)~Horas de viagem por veículo-mês, por tipo de uso, relativas à mediana. (b)~Fração acumulada dos acidentes e das horas contra a fração acumulada de veículos, ordenados do maior para o menor; faixa: 95\% das curvas simuladas sob Poisson com taxa constante por tipo de uso. (c)~Distribuição dos acidentes por hora do dia; faixas: turnos de 8 horas. (d)~Relatos de acidente, ordens de manutenção por dano e colisões detectadas pela telemetria, por mês, relativos à média de cada série.}%
  {Preliminary description of the operational data (relative values only). (a)~Trip hours per vehicle-month, by use type, relative to the median. (b)~Cumulative share of accidents and hours against the cumulative share of vehicles, ordered from largest to smallest; band: 95\% of curves simulated under Poisson with a constant rate per use type. (c)~Distribution of accidents by hour of day; bands: 8-hour shifts. (d)~Accident reports, damage repair orders and telematics-detected collisions per month, relative to each series' mean.}}
\label{fig:desc-pp}
\end{figure}

\begin{nota}[\tr{Privacidade}{Privacy}]
\tr%
  {Nenhum dado identifica pessoas. A análise usa veículo, unidade e mês; identificadores de motorista, nomes, imagens de câmera e trajetos de GPS não são usados. Os dados não saem do ambiente da empresa; apenas tabelas agregadas entram no repositório da pesquisa e na dissertação, que não nomeará a empresa, seus clientes nem suas unidades.}%
  {No data identifies people. The analysis uses vehicle, site and month; driver identifiers, names, camera footage and GPS tracks are not used. Data stays within the company's environment; only aggregated tables enter the research repository and the dissertation, which will not name the company, its clients or its sites.}
\end{nota}
